\documentclass[12pt,aspectratio=16:9]{beamer}
\usepackage[utf8]{inputenc}
\usepackage[T1]{fontenc}
\usepackage[portuguese]{babel}
\usepackage{graphicx}
\usepackage{booktabs}
\usepackage{tikz}
\usetikzlibrary{arrows.meta,positioning,shapes.geometric,calc}
\usepackage{listings}
\usepackage{xcolor}
\usepackage{multimedia}
\usepackage{hyperref}

% Color definitions
\definecolor{primary}{RGB}{18,56,102}
\definecolor{secondary}{RGB}{28,120,86}
\definecolor{accent}{RGB}{220,53,69}
\definecolor{lightblue}{RGB}{236,242,250}
\definecolor{lightgreen}{RGB}{235,247,241}

% Beamer theme
\usetheme{Madrid}
\usecolortheme[named=primary]{structure}
\setbeamercolor{title}{fg=white,bg=primary}
\setbeamercolor{frametitle}{fg=white,bg=primary}
\setbeamercolor{block title}{fg=white,bg=secondary}
\setbeamercolor{block body}{bg=lightgreen}

% Title page
\title{Painel IoT Distribuído com Monitoramento em Tempo Real}
\subtitle{Infraestrutura Heterogênea: Raspberry Pi, Linux Mint e ESP32/ESP8266}
\author{Leonardo Cardoso de Moura \\ \textit{Matrícula: 32021BSI004}}
\date{08 de maio de 2026}
\institute{Disciplina: Sistemas Distribuídos}

% Code style
\lstset{
    language=C,
    basicstyle=\ttfamily\small,
    backgroundcolor=\color{lightblue},
    breaklines=true,
    frame=single,
    rulecolor=\color{primary},
    showstringspaces=false,
    keywordstyle=\color{primary},
    commentstyle=\color{gray},
    stringstyle=\color{secondary}
}

\begin{document}

% Title slide
\frame{\titlepage}

% Outline
\begin{frame}{Roteiro}
    \tableofcontents
\end{frame}

% ==========================
\section{Introdução}
% ==========================

\begin{frame}{O Micro-mundo}
    \begin{center}
        \textbf{Sistema IoT Distribuído para Monitoramento Ambiental em Tempo Real}
    \end{center}
    
    \vspace{1cm}
    
    \begin{columns}
        \column{0.5\textwidth}
        \begin{block}{Características}
            \begin{itemize}
                \item Coleta contínua de dados
                \item Múltiplos nós geograficamente distribuídos
                \item Sincronização centralizada
                \item Dashboard web
                \item Displays locais OLED
            \end{itemize}
        \end{block}
        
        \column{0.5\textwidth}
        \begin{block}{Sensores}
            \begin{itemize}
                \item Temperatura
                \item Umidade
                \item Pressão
                \item Força de sinal WiFi (RSSI)
                \item Qualidade de conexão
            \end{itemize}
        \end{block}
    \end{columns}
\end{frame}

\begin{frame}{O que o Sistema Faz}
    \begin{enumerate}[<+->]
        \item Coleta contínua de leituras de sensores (temperatura, umidade, pressão)
        \item Monitoramento de força e qualidade de sinal WiFi para diagnóstico
        \item Exibição em tempo real em telas OLED com animações responsivas
        \item Controle interativo via sensor de toque capacitivo
        \item Sincronização com servidor central através de API HTTP
        \item Persistência de séries temporais em PostgreSQL
        \item Dashboard web para visualização e controle remoto
    \end{enumerate}
\end{frame}

% ==========================
\section{Arquitetura do Sistema}
% ==========================

\begin{frame}{Arquitetura Geral}
    \centering
    \begin{tikzpicture}[
        node distance=2.5cm and 2.5cm,
        esp/.style={rectangle, rounded corners=3mm, draw=secondary, very thick, 
                   minimum width=2.8cm, minimum height=0.8cm, align=center, fill=lightgreen!70},
        server/.style={rectangle, rounded corners=3mm, draw=primary, very thick, 
                      minimum width=2.8cm, minimum height=0.8cm, align=center, fill=lightblue!70},
        arrow/.style={-{Latex[length=2.5mm]}, very thick, draw=gray}
    ]

    \node[esp] (esp1) {ESP32 (Sala)};
    \node[esp, below right=2cm and 1.5cm of esp1] (esp2) {ESP8266 (Varanda)};
    
    \node[server, below=4cm of esp1] (nginx) {Nginx};
    \node[server, right=2.2cm of nginx] (api) {API Node.js};
    \node[server, right=2.2cm of api] (db) {PostgreSQL};

    \draw[arrow] (esp1) -- (nginx);
    \draw[arrow] (esp2) -- (nginx);
    \draw[arrow] (nginx) -- (api);
    \draw[arrow] (api) -- (db);
    
    \node[font=\small, above left=0.2cm of nginx] {WiFi POST};
    \node[font=\small, above=0.2cm of api] {HTTP};
    \node[font=\small, above=0.2cm of db] {TCP};
    \end{tikzpicture}
\end{frame}

\begin{frame}{Infraestrutura Física}
    \begin{columns}
        \column{0.5\textwidth}
        \begin{block}{\textbf{Raspberry Pi 3B}}
            \begin{itemize}
                \item PostgreSQL
                \item Banco de dados persistente
                \item IP fixo: 192.168.1.10
                \item Conectividade: Ethernet
            \end{itemize}
        \end{block}
        
        \column{0.5\textwidth}
        \begin{block}{\textbf{Notebook (Linux Mint)}}
            \begin{itemize}
                \item Node.js API
                \item Nginx proxy reverso
                \item IP fixo: 192.168.1.11
                \item Conectividade: WiFi
            \end{itemize}
        \end{block}
    \end{columns}
    
    \vspace{0.5cm}
    
    \begin{block}{\textbf{ESP32/ESP8266 (Distribuídos)}}
        \begin{itemize}
            \item Múltiplos nós em diferentes locais
            \item Coleta de dados via sensores
            \item WiFi para comunicação
            \item Displays OLED locais
        \end{itemize}
    \end{block}
\end{frame}

\begin{frame}{Fluxo de Comunicação}
    \begin{center}
        ESP32/ESP8266 $\rightarrow$ \textbf{[WiFi]} $\rightarrow$ Nginx $\rightarrow$ API Node.js $\rightarrow$ PostgreSQL
    \end{center}
    
    \vspace{1cm}
    
    \begin{block}{Protocolo de Dados}
        \textbf{HTTP POST com payload JSON}
    \end{block}
    
    \vspace{0.5cm}
    
    \begin{columns}
        \column{0.5\textwidth}
        \begin{itemize}
            \item Frequência: 30-60s
            \item Buffer local: 10 leituras
            \item Auto-recover WiFi
            \item Backoff exponencial
        \end{itemize}
        
        \column{0.5\textwidth}
        \begin{itemize}
            \item Rate limiting
            \item Compressão gzip
            \item Logging de requisições
            \item CORS habilitado
        \end{itemize}
    \end{columns}
\end{frame}

% ==========================
\section{Sensores e Dispositivos}
% ==========================

\begin{frame}{Sensores Instalados}
    \begin{columns}
        \column{0.5\textwidth}
        \begin{block}{\textbf{BMP280}}
            \begin{itemize}
                \item Temperatura: -40 a +85°C
                \item Pressão: 300-1100 hPa
                \item Protocolo: I2C/SPI
                \item Precisão: ±1°C
            \end{itemize}
        \end{block}
        
        \column{0.5\textwidth}
        \begin{block}{\textbf{DHT11}}
            \begin{itemize}
                \item Temperatura: 0-50°C
                \item Umidade: 20-80\%
                \item Protocolo: Single-wire
                \item Precisão: ±2°C, ±5\%
            \end{itemize}
        \end{block}
    \end{columns}
    
    \vspace{0.8cm}
    
    \begin{block}{\textbf{RSSI (Received Signal Strength Indicator)}}
        Monitoramento de qualidade de sinal WiFi:
        \begin{itemize}
            \item Diagnóstico real de cobertura
            \item Detecção de interferência
            \item Planejamento de repetidores
        \end{itemize}
    \end{block}
\end{frame}

\begin{frame}{Displays e Interatividade}
    \begin{columns}
        \column{0.6\textwidth}
        \begin{block}{\textbf{Display OLED SSD1306}}
            \begin{itemize}
                \item 128×64 pixels
                \item Protocolo I2C
                \item Taxa: 1-2 Hz
                \item Conteúdo:
                \begin{itemize}
                    \item Temperatura em tempo real
                    \item Força de sinal WiFi
                    \item Animações de rosto
                    \item Status do sistema
                \end{itemize}
            \end{itemize}
        \end{block}
        
        \column{0.4\textwidth}
        \begin{block}{\textbf{Sensor Toque TTP223}}
            \begin{itemize}
                \item Controle capacitivo
                \item 5 expressões
                \item Debounce 300ms
                \item Ciclagem ao tocar
            \end{itemize}
            
            \vspace{0.5cm}
            
            \centering
            \small
            😊 😢 😲 🤔 ❤️
        \end{block}
    \end{columns}
\end{frame}

% ==========================
\section{Banco de Dados}
% ==========================

\begin{frame}{Schema PostgreSQL}
    \begin{block}{\textbf{Tabela: telemetry} (TimescaleDB)}
        \begin{itemize}
            \item \texttt{time} (TIMESTAMP com timezone)
            \item \texttt{device\_id} (chave de particionamento)
            \item \texttt{temperature\_c, humidity\_percent, pressure\_hpa}
            \item \texttt{rssi\_dbm, battery\_voltage}
            \item Retenção: 90 dias
        \end{itemize}
    \end{block}
    
    \begin{block}{\textbf{Tabela: devices}}
        \begin{itemize}
            \item ID, nome, localização, status
            \item Last seen, firmware version
            \item IP address (último visto)
        \end{itemize}
    \end{block}
    
    \begin{block}{\textbf{Tabela: alerts}}
        Tipo, severidade, timestamp criação/resolução
    \end{block}
\end{frame}

\begin{frame}{Payload JSON Exemplo}
    \lstset{language=JSON, basicstyle=\ttfamily\tiny}
    \begin{lstlisting}
{
  "device_id": "ESP32_SALA_001",
  "timestamp": 1715179200,
  "sensor_data": {
    "temperature_c": 24.5,
    "humidity_percent": 65.2,
    "pressure_hpa": 1013.5,
    "rssi_dbm": -45,
    "battery_voltage": 3.3
  },
  "display_state": {
    "current_expression": "happy",
    "screen_active": true
  },
  "system_health": {
    "wifi_reconnects": 2,
    "api_failures": 0
  }
}
    \end{lstlisting}
\end{frame}

% ==========================
\section{Conceitos de Sistemas Distribuídos}
% ==========================

\begin{frame}{Objetivos de Sistemas Distribuídos}
    \begin{block}{1. \textbf{Escalabilidade}}
        Novos nós podem ser adicionados sem redesenho. Nginx permite escalabilidade horizontal.
    \end{block}
    
    \begin{block}{2. \textbf{Disponibilidade}}
        Buffer local em ESP, Nginx e API redundantes, PostgreSQL com snapshots automáticos.
    \end{block}
    
    \begin{block}{3. \textbf{Tolerância a Falhas}}
        Reconexão automática, retry com backoff, health check, fila de mensagens locais.
    \end{block}
    
    \begin{block}{4. \textbf{Compartilhamento de Recursos}}
        Nginx como ponto centralizado, PostgreSQL como estado único, transações ACID.
    \end{block}
\end{frame}

\begin{frame}{Desafios Reais do Sistema}
    \begin{columns}
        \column{0.5\textwidth}
        \begin{block}{\textbf{Infraestrutura}}
            \begin{itemize}
                \item IP dinâmico do provedor
                \item Falhas WiFi frequentes
                \item Latência variável
                \item Nós geograficamente dispersos
            \end{itemize}
        \end{block}
        
        \column{0.5\textwidth}
        \begin{block}{\textbf{Soluções}}
            \begin{itemize}
                \item DHCP estático no roteador
                \item Buffer local nos nós
                \item Backoff exponencial
                \item Sincronização periódica
            \end{itemize}
        \end{block}
    \end{columns}
\end{frame}

% ==========================
\section{Roadmap Futuro}
% ==========================

\begin{frame}{Próximas Etapas de Desenvolvimento}
    \begin{itemize}[<+->]
        \item \textbf{Etapa 2}: Threads em Node.js, processamento paralelo
        \item \textbf{Etapa 3}: UDP para comunicação de baixa latência
        \item \textbf{Etapa 4}: RPC entre servidores
        \item \textbf{Etapa 5}: Message Bus (RabbitMQ/MQTT)
        \item \textbf{Etapa 6}: Service Discovery (mDNS/Consul)
        \item \textbf{Etapa 7}: Mutual exclusion, eleição de líder
        \item \textbf{Etapa 8}: Replicação com consistência eventual
        \item \textbf{Etapa 9}: Circuit breaker, resiliência
        \item \textbf{Etapa 10}: TLS, autenticação mútua, JWT
    \end{itemize}
\end{frame}

% ==========================
\section{Conclusão}
% ==========================

\begin{frame}{Conclusão}
    \begin{block}{\textbf{Um Sistema Distribuído Real}}
        \begin{itemize}
            \item Infraestrutura genuinamente prática e doméstica
            \item Múltiplos nós heterogêneos (Raspberry, Notebook, ESP32/ESP8266)
            \item Rede WiFi com falhas esperadas
            \item Persistência centralizada
            \item Desafios reais de sincronização e resiliência
        \end{itemize}
    \end{block}
    
    \vspace{1cm}
    
    \begin{center}
        \Large \textbf{Fundação Sólida para Conceitos Avançados}
    \end{center}
\end{frame}

\begin{frame}
    \centering
    \vspace{2cm}
    {\huge \textbf{Obrigado!}}
    
    \vspace{1cm}
    
    {\Large Leonardo Cardoso de Moura}\\
    {\small Matrícula: 32021BSI004}\\
    {\small Disciplina: Sistemas Distribuídos}
    
    \vspace{2cm}
    
    {\small 08 de maio de 2026}
\end{frame}

\end{document}
